% ==============================================================================
% SECCIÓN 6: ANTECEDENTES DEL EQUIPO
% ==============================================================================
\newpage
\section{Antecedentes del Equipo de Trabajo}

La propuesta identifica al equipo CLENCH Software Development. Para presentar su capacidad de ejecución con evidencias, cada integrante completará conocimientos relevantes, experiencia previa y una referencia verificable. No se atribuyen tecnologías ni experiencia que todavía no hayan sido confirmadas por sus integrantes.

\begin{center}
\small
\begin{tabularx}{\textwidth}{|p{3.5cm}|X|p{3.4cm}|}
\hline
\rowcolor{tableHeader}
Integrante & Experiencia y conocimientos relacionados con el proyecto & Evidencia / referencia \\
\hline
Miguel Ángel Luque Gómez & Full-Stack, desarrollo de aplicaciones web y móviles, conocimientos en UX/UI & SIBW, DIU\\
\hline
Yeray Rodríguez Navas & Full-stack, creación de aplicaciones web y móviles, experiencia en UX/UI & DS, SIBW, DIU \\
\hline
Julián Carrión Tovar & Full-stack, creación de plataformas web y aplicaciones móviles, experiencia en UX/UI. & \href{https://juliancarrion.dev/}{Portafolio} \\
\hline
Pablo Hernández Ibáñez & Full-Stack, creación de aplicaciones web y móviles (Flutter), experiencia en UX/UI. Experiencia en entorno de interacción y escenas 3D. & DS, SIBW, DIU\\
\hline
Pablo de la Torre Roldán & Full Stack development para el desarrollo de páginas web (Flutter, PHP), experiencia en UX/UI, experiencia en desarrollo de aplicaciones móviles con Flutter & DS, SIBW, DIU,  \href{https://wpd.ugr.es/~sig.proc.audiology/}{SigProcAudiology}\\
\hline
José Rodríguez Fernández & Full-stack, creación de aplicaciones web, móvil y experto en usabilidad & DS/SIBW \\
\hline
Amine Azzammouri & Full-stack, aplicaciones web y experiencia en UX/UI & DIU, SIBW, PW \\
\hline
\end{tabularx}
\end{center}

\subsection{Capacidades que debe cubrir el equipo}
Al completar la tabla, se comprobará que el equipo dispone de responsables o apoyo para: desarrollo de la aplicación, backend y persistencia, diseño accesible de interfaces, pruebas y documentación, y coordinación con el centro.